\begin{recipe}{Tomaten-pestorisotto}
\kicker[Hoofdgerecht,Vegetarisch,Italiaans-Hollands]{Hoofdgerecht · vegetarisch · Italiaans-Hollands}
\index[register]{Tomaten-pestorisotto}
\index[register]{Risotto, tomaten-pesto-}

\meta{35 minuten}

\lettrine{R}{omige} risotto met zoete, in de oven geroosterde trostomaatjes en
een lepel verse pesto. Bij ons zorgt dit gerecht steevast voor een lege pan,
en voor de vraag om nog een keer.

\blockrule{Ingrediënten}
\begin{ingredients}
  \ing{4}{vleestomaten, gewassen en in blokjes}
  \ing{250 g}{trostomaatjes (aan de tak)}\index[register]{Trostomaatjes, geroosterd}
  \ing{3 el}{verse groene pesto}\index[register]{Pesto, verse}
  \ing{170 g}{risottorijst}\index[register]{Risotto, tomaten-pesto-}
  \ing{600 ml}{groentebouillon}\index[register]{Bouillon, groente-}
  \ing{1}{grote ui, gesnipperd}\index[register]{Ui}
  \ing{2}{knoflookteentjes, fijngehakt}\index[register]{Knoflook}
  \ing{150 ml}{droge witte wijn}\index[register]{Witte wijn}
  \ing{75 g}{parmezaan, geraspt}\index[register]{Parmezaanse kaas}
  \ing{30 g}{ongezouten roomboter}
  \ing{2 el}{olijfolie}
  \ing{1 bos}{verse basilicum}\index[register]{Basilicum}
\end{ingredients}

\blockrule{Bereiding}
\tip{Voeg de bouillon beetje bij beetje toe en blijf roeren. Dat geduld maakt
de risotto pas écht romig.}
\begin{steps}
  \item Houd de groentebouillon op een laag vuur warm; je voegt hem straks lepel
        voor lepel toe.
  \item Verwarm ondertussen de oven voor op 200\,°C (boven- en onderwarmte), zodat 
        hij op temperatuur is wanneer de tomaatjes de oven in gaan.
  \item Fruit de ui en knoflook zachtjes in olijfolie in een hapjespan. Bak de
        risottorijst kort mee tot de korrels glazig zien. 
  \item Was de trostomaatjes en zet ze in de oven, op een bakplaat met een scheutje olijfolie.
  \item Blus af met de witte wijn en laat die volledig verdampen.
  \item Voeg de warme bouillon geleidelijk toe, telkens een soeplepel, en blijf
        roeren tot de rijst beetgaar en romig is.
  \item Roer de verse tomaten in twee fasen door de risotto: de eerste geeft de
        saus zoetheid mee, de tweede houdt wat frisse bite.
  \item De trostomaatjes staan intussen zo'n 20 tot 25 minuten in de oven;
        haal ze eruit zodra ze zacht en zoet zijn.
  \item Roer van het vuur af de boter en Parmezaan erdoor. Laat de risotto een
        paar minuten rusten, dan wordt hij nog romiger en glanzender.
  \item Schep de basilicum erdoor en serveer met een lepel verse pesto en de
        geroosterde trostomaatjes.
\end{steps}

\vspace{1em}{\footnotesize\itshape\color{terracotta}Met dank aan Miljuschka
Witzenhausen, naar haar tomaten-basilicumrisotto\cite{witzenhausen-risotto}.}

\heroimagefade{images/risotto}
\end{recipe}
